\section{Visualisierungen}
Zur Beantwortung der zentralen Fragestellung muss der Nutzer die im Folgenden vorgestellten Anwendungsaufgaben bearbeiten. Daraus werden Anforderungen an eine Visualisierung abgeleitet, die den Nutzer bei der Bearbeitung unterstützen sollen. Im Anschluss werden die gewählten Visualisierungen vorgestellt und die getroffenen Designentscheidungen diskutiert.
\subsection{Analyse der Anwendungsaufgaben}\label{subsec:aufgabe}
% [ ] Analysieren Sie die konkreten Anwendungsaufgaben, die zur Lösung des Zielproblems durch die Anwender bearbeitet werden müssen. --Neue Perspektive bieten, um zu erkennen, dass der Medaillenspiegel nicht ausreichend ist--
% [ ] Was müssen Anwender angucken, um ihre Aufgaben zu lösen
% --> Nicht die Visualisierung beschreiben. Die Visualisierung soll das Ergebnis der Diskussion sein.

% [ ] Welche sinnvollen mentale Modelle helfen den Personen bei der Bearbeitung. 
% -> Mentale Modelle: geometrische Objekte, Zeitreihen, Zustandsänderungen
% [ ] Sind diese mentalen Modelle für sie notwendig, um die Aufgaben lösen zu können?
% --> Gehen Sie bei Ihrer Argumentation von den Anwendungsaufgaben aus und kommen Sie dann zu den mentalen Modellen, deren Aufbau durch Visualisierungen unterstützt wird.
% [ ] Welche Visualisierungen helfen den Personen, die die Software verwenden, sinnvolle mentale Modelle aufzubauen. 

Die Anwender sollen mithilfe dieser Visual-Analytics-Anwendung die Frage beantworten können, welche Nation die sportlich erfolgreichste bei den Olympischen Sommerspielen 2024 war. Dabei sollen sie erkennen, dass diese Frage nicht allein mit dem Medaillenspiegel beantwortet werden kann, da dieser nur die Anzahl gewonnener Medaillen darstellt. Stattdessen sollen die Anwender die Ergebnisse detaillierter und im Kontext demografischer und wirtschaftlicher Einflussfaktoren bewerten. Dazu muss die Visual-Analytics-Anwendung ihnen eine neue Perspektive auf die Ergebnisse der Olympischen Sommerspiele 2024 bieten.

\paragraph{Analyse Eins}
% Analysieren Sie die konkrete Anwendungsaufgabe, die zur Lösung des Zielproblems bearbeitet werden muss
Die Nutzer beginnen die Analyse bei dem bereits bekannten Medaillenspiegel und sollen sich in einem ersten Schritt über die Herkunft der gewonnenen Medaillen informieren. In dem Medaillenspiegel können sie lediglich die Anzahl gewonnener Medaillen ablesen. Daher ist die erste Anwendungsaufgabe, die Verteilung der Medaillen nach Sportarten und Disziplinen zu analysieren. Bei den Olympischen Sommerspielen 2024 fanden insgesamt 329 Wettkämpfe in 32 Sportarten statt. Die meisten Nationen sind aber nur in wenigen Sportarten erfolgreich. Südostasiatischen Ländern sind häufig nur in den Kampfsportarten erfolgreich. Eine sportlich erfolgreiche Nation kann Medaillengewinner aber in vielen verschiedenen Sportarten beglückwünschen.

%Was müssen Anwender angucken, um ihre Aufgaben zu lösen
Um sportlich vielseitige Nationen zu identifizieren, müssen die Anwender erkennen können, in welchen Sportarten und Disziplinen die Medaillen gewonnen wurden und welchen prozentualen Anteil dies umfasst. Die Vereinigten Staaten haben zwar die meisten Medaillen gewonnen, jedoch stammen über 50\% davon aus den Sportarten \textit{Leichtathletik} und \textit{Wassersport}. Diese Verteilung soll für die Anwender leicht erkennbar sein.

%Welche sinnvollen mentale Modelle helfen den Personen
Typischerweise werden solche Daten als Anteile eines Kreises modelliert. Auch wenn die einfache Auflistung der gewonnenen Medaillen pro Sportart zur Beantwortung der Frage ausreichen würde, kann durch das mentale Modell des Kreises die prozentuale Größe leichter erfasst werden, da 50\% auch der Hälfte des Kreises entspricht.

\paragraph{Analyse Zwei}
% Analysieren Sie die konkrete Anwendungsaufgabe, die zur Lösung des Zielproblems bearbeitet werden muss
Die in dem Medaillenspiegel abgebildeten sportlichen Erfolge resultieren aus den demografischen und wirtschaftlichen Entwicklungen dieser Nation. Finanzielle Förderungen sind die Grundlage für sportliche Erfolge, die sich aber nicht jedes Land leisten kann. Jedoch wird dieser Zusammenhang in der rein quantitativen Analyse des Medaillenspiegels nicht beachtet. Aus diesem Grund dominieren seit jeher die Industrienationen die vorderen Plätze in diesem Ranking. Zur Beantwortung der Frage nach der sportlich erfolgreichsten Nation soll der Nutzer die demografischen und wirtschaftlichen Einflussfaktoren in Beziehung zu den sportlichen Erfolgen eines Landes stellen. Dadurch findet in gewisser Weise eine Normalisierung des Medaillenspiegels statt, wobei vor allem die wirtschaftlichen Vorteile der Industrienationen gegenüber Entwicklungsländern relativiert werden. Für eine genauere Untersuchung werden umfangreichere Ergebnisse und komplexere Einflussfaktoren aus der Soziologie und den Wirtschaftswissenschaften benötigt, welche an dieser Stelle nicht weiter vertieft werden.

%Was müssen Anwender angucken, um ihre Aufgaben zu lösen
Nachdem die Medaillen in Beziehung zu demografischen und wirtschaftlichen Faktoren gesetzt wurden, soll der Nutzer die neuen Länderrankings mit dem Medaillenspiegel vergleichen, um sportlich erfolgreiche Nationen, relativ zu den qualitativen Faktoren, erkennen zu können. Der Anwender soll also die Platzierung eines Landes in den verschiedenen Länderrankings vergleichen, um dadurch auf die unterschiedliche Bewertung der sportlichen Erfolge einer Nation im Kontext äußerer Einflussfaktoren aufmerksam zu werden.

%Welche sinnvollen mentale Modelle helfen den Personen
Rankings werden typischerweise tabellarisch dargestellt, wobei sich der Bestplatzierte an oberster Stelle befindet. Für diese Anwendungsaufgabe ist dieses mentale Modell nicht notwendig, da in erster Linie Veränderungen bei den Platzierungen in den verschiedenen Rankings betrachtet werden sollen. Nichtsdestotrotz ist eine gesamte Darstellung des Rankings nützlich, um einen Überblick zu erhalten. Dabei dient das mentale Modell eines tabellarischen Rankings den Anwendern zur schnellen Orientierung.

\paragraph{Analyse Drei}
% Analysieren Sie die konkrete Anwendungsaufgabe, die zur Lösung des Zielproblems bearbeitet werden muss
Während der Medaillenspiegel nur die Ergebnisse einer Olympiade zusammenfasst, können sportlich erfolgreiche Nationen regelmäßig neue Medaillengewinner hervorbringen, da sie eine gute Talententwicklung und sportliche Förderungen haben. Die sportlichen Erfolge von Entwicklungsländern beruhen dagegen häufig auf einzelnen Ausnahmeathleten, welche zum Teil nicht einmal im eigenen Land trainieren. Außerdem muss bei dieser Analyse beachtet werden, dass die gezielte Talententwicklung ein Prozess ist, deren Erfolge erst einige Jahre später zu erkennen ist. Jedoch kann dieser Prozess durch immer häufigere Medaillengewinner beobachtet werden. Aus diesem Grund soll der Nutzer als drittes die erzielten Ergebnisse der Olympischen Sommerspiele 2024 mit vorhergehenden Jahren vergleichen, um die sportlich erfolgreichen Nationen an regelmäßigen Medaillengewinnern bzw. steigenden Entwicklungen im Medaillenspiegel erkennen können.

%Was müssen Anwender angucken, um ihre Aufgaben zu lösen
Für diese Analyse muss der Nutzer die gewonnenen Medaillen der Olympischen Sommerspiele 2024, als auch die aller vergangenen Olympiaden angucken und miteinander vergleichen. Anhand dieser Daten soll der Nutzer Nationen mit regelmäßigen Medaillengewinnern, sowie sportlich aufstrebende Nationen identifizieren.

%Welche sinnvollen mentale Modelle helfen den Personen
Die historische Entwicklung des Medaillenspiegels einer Nation interpretieren viele Anwender vermutlich als Zeitstrahl. Auf diesem werden in zeitlich fortschreitenden Abständen die Anzahl gewonnener Medaillen dargestellt. Dieses mentale Modell kann auch von Nutzern ohne Vorwissen intuitiv interpretiert werden. Für die Beantwortung der Fragestellung müssen die Anwender aber auch die Unterschiede zwischen den einzelnen Jahren analysieren können. Dabei sind die konkreten Werte nicht so wichtig, wie die Größe der Veränderung. Ein mentales Modell, das bei dieser Aufgabe hilft, sind Farben, durch die große Änderungen mehr hervorgehoben werden können als kleine Änderungen.

\subsection{Anforderungen an die Visualisierungen}
% Leiten Sie Anforderungen an das Design der Visualisierungen ab, die sich durch Ihre Analyse des Zielproblems ergeben.
% Trivial: soll verständlich sein
% Besser: Welche Kompromisse, was besonders gut zu dekodieren sein soll
\paragraph{Visualisierung Eins}
Die Anwender sollen die Verteilung der Medaillen einer Nation nach Sportarten und Disziplinen analysieren. Als mentales Modell wurden Anteile eines Kreises identifiziert, da so die prozentualen Größen im richtigen Verhältnis dargestellt werden können. Diese Eigenschaft sollte in der Visualisierung ebenfalls umgesetzt werden, da es die Bearbeitung der Anwendungsaufgabe erleichtert. Weiterhin unterstützen Farben bei der Identifikation der Anteile. Die Visualisierung sollte die verschiedenen Anteile durch unterschiedliche Farben hervorheben, sodass große Flächen gleicher Farbe (viele Medaillen in der gleichen Sportart) schnell erkannt werden können.

Außerdem sind viele Sportarten Oberkategorien zugeordnet und in Disziplinen unterteilt. Diese Hierarchie sollte durch die Visualisierung ebenfalls hervorgehoben und hierarchisch zusammengehörige Anteile zusammengefasst werden.

Bei den Vereinigten Staaten muss die Verteilung von 126 Medaillen dargestellt werden. Die Darstellung aller dazugehörigen Informationen würde die Visualisierung überladen und unübersichtlich machen. Zur Bearbeitung der Anwendungsaufgabe sind die konkreten Namen der Sportarten und Disziplinen nur zweitrangig und können temporär (z.B. per Mouse-Hover) eingeblendet werden. 

\paragraph{Visualisierung Zwei}
Der Medaillenspiegel wird typischerweise als Tabelle, mit der erfolgreichsten Nation in der obersten Zeile, dargestellt. Damit die Anwender diese Vorstellung beibehalten können, sollten der Vergleich mit den qualitativen Einflussfaktoren ebenfalls in einem Länderranking resultieren, bei dem sich der Sieger, also die sportlich erfolgreichste Nation, an oberster Stelle befindet. 

Gleichzeitig sollen die Anwender aber auch Unterschiede zwischen den Platzierungen eines Landes in den verschiedenen Länderrankings erkennen können. Das heißt, die Länderrankings müssen visuell schnell vergleichbar sein. Dies ist bei der Darstellung aller 204 teilnehmenden Länder als Tabelle nicht möglich. Ein Kompromiss ist die Darstellung der Platzierung in dem jeweiligen Länderranking, ohne der konkreten Werte. Die Vergleichbarkeit wäre weiterhin möglich. 

\paragraph{Visualisierung Drei}
Für die Analyse des historischen Medaillenspiegels müssen die Anwender die gewonnenen Medaillen der vergangenen Olympiaden miteinander vergleichen. Seit Beginn der Olympischen Sommerspiele der Neuzeit 1896 gab es 33 Olympiaden, wobei diese in den Jahren 1916, 1940, 1944 wegen Weltkriegen nicht stattfinden konnten. Die Spiele finden nur alle 4 Jahre statt, sodass die Darstellung mit dem mentalen Modell des Zeitstrahls viele Leerräume hätte. Die Visualisierung sollte die Daten daher möglichst effektiv darstellen und nur die Jahre darstellen, in denen auch eine Olympiade stattgefunden hat.

Bei dieser Aufgabe sind die Differenzen zwischen den Werten (gewonnene Medaillen) gering, da die meisten Länder in der Regel nicht mehr als 10 Medaillen gewinnen. Für die Analyse der Differenzen wurden das mentale Modell von Farbänderungen identifiziert. Die Änderungen einer Farbe sollten für kleine Werte groß sein, um die Unterschiede visuell deutlich hervorzuheben. Da die Zielgruppe unterschiedliche kulturelle Erfahrungen mitbringt sollte außerdem sichergestellt werden, dass alle Anwender die Bedeutung der Farben gleich interpretieren.

\subsection{Präsentation der Visualisierungen}
Aus den formulierten Anwendungsaufgaben wurden Anforderungen an eine Visualisierung formuliert, welche den Anwender bei der Bearbeitung unterstützen sollen. Im Folgenden werden die daraus abgeleiteten Visualisierungen vorgestellt und deren Vorteile bei der Bearbeitung der Anwendungsaufgabe diskutiert. 
% [ ] Präsentieren sie die visuelle Abbildungen und Kodierungen der Daten und Interaktionsmöglichkeiten. 
% [ ] Begründen sie, warum und wie gut ihre Designentscheidungen die erstellten Anforderungen erfüllen. 
% [ ] Begründen sie, warum die gewählte visuelle Kodierung der Daten für das zulösenden Problem passend ist.
% --> Typische Argumente würden hier auf Wahrnehmungsprinzipien und Theorie über Informationsvisualisierung verweisen. 
% --> Die besten Begründungen diskutieren explizit die konkrete Auswahl der Visualisierungen im Kontext von mehreren verschiedenen Alternativen. 
% --> Machen sie hier nicht den Fehler, einfach nur Visualisierung aus den vorgegebenen Bereichen zu diskutieren, weil das in der Regel nicht sinnvoll ist.
% --> Wenn sie sich für einen Scatterplot entschieden haben, ist ein Zeitreihendiagramm in der Regel keine Alternative.
% --> Diskutieren sie also nicht einfach Zeitreihendiagramme, weil sie in den Anforderungenen an das Projekt neben Scatterplots stehen, sondern suchen sie nach echten alternativen Visualisierungen, die zum Aufbau eines vergleichbaren mentalen Modells führen. 
% [ ] Diskutieren sie die Expressivität und die Effektivität der einzelnen Visualisierungen.
\subsubsection{Visualisierung Eins}
% Einleitung
Diese Visualisierung soll die Bearbeitung der Anwendungsaufgabe auf zwei verschiedenen Arten unterstützen. Zum einen sollen die prozentualen Anteile gewonnener Medaillen pro Sportart dargestellt werden. Das mentale Modell von Anteilen eines Kreises kann zum Beispiel durch ein Kreisdiagramm umgesetzt werden. Zum anderen soll die Zuordnung einer Sportart zu einer Kategorie und die Untergliederung in verschiedene Disziplinen dargestellt werden. Das SunBurst-Diagramm ist eine spezielle Form des Kreisdiagramms, welches neben der Darstellung der prozentualen Anteile auch diese Hierarchie abbilden kann.

% Präsentation der visuellen Abbildung und Kodierung der Daten
Jede gewonnene Medaille wird zunächst durch einen Anteil entsprechend der insgesamt gewonnenen Medaillen dargestellt. Anschließend werden alle hierarchisch zusammengehörigen Medaillen vereinigt und auf einer neuen Hierarchieebene dargestellt. Gegebenenfalls werden die verschiedenen Sportarten noch zu einer gemeinsamen Kategorie vereinigt (siehe \textit{Wassersport}) und durch eine dritte Hierarchieebene dargestellt. Schließlich werden die dargestellten Anteile durch Farben hervorgehoben, wobei darauf geachtet wurde, dass jede Sportart eine andere Farbe erhält.

% Begründung, warum und wie gut Desingentscheidung die Anforderungen erfüllt
Die Visualisierungstechnik SunBurst bildet die prozentualen Anteile im richtigen Verhältnis ab, sodass 100\% auch durch einen vollständigen Kreis dargestellt werden. Die prozentuale Größe ist dadurch leicht zu erfassen. Die Anteile im Diagramm wurden mit unterschiedlichen Farben gezeichnet, sodass diese neben der Größe zusätzlich auch durch die Farbe visuell unterschieden werden können. Die dritte Anforderung zur Darstellung der Untergliederung wird durch die verschiedenen Hierarchieebenen erfüllt. Damit konnten alle Anforderungen an die Visualisierung umgesetzt werden.

% Begründung, warum gewählte visuelle Kodierung der Daten für Anwendungsaufgabe passend ist (z.B. Wahrnehmungsprinzip, Theorie über InfoVis)
Die Visualisierung soll den Anwender dabei unterstützen, die Anteile an einem Ganzen zu analysieren. Das mentale Modell von Anteilen eines Kreises ist für die meisten Anwender die intuitive Wahl. Anteile können aber auch durch andere geometrische Objekte, wie zum Beispiel Rechtecke, dargestellt werden. Diese werden in Balkendiagrammen gestapelt und in TreeMaps ineinander geschachtelt verwendet. Durch das Ineinander-Schachteln können mit TreeMaps auch Hierarchien, wie die Untergliederung von Sportarten in Disziplinen, dargestellt werden. Die TreeMaps würden damit die Anforderungen an die Visualisierung genauso erfüllen. Wir haben uns trotzdem für das SunBurst-Diagramm entschieden, da dieses verbreiteter ist und unsere Anwendung auch von einer Zielgruppe ohne Vorwissen genutzt werden soll. Für die Interpretation von TreeMaps müssen die Anwender die Bedeutung von ineinander-geschachtelten Rechtecken kennen, um die Hierarchie zu erfassen. Dies kann von ungeübten Nutzer nicht erwartet werden. In einem SunBurst-Diagramm wird die Hierarchie durch konzentrische Ringe von innen nach außen dargestellt, wodurch eine gewisse Ähnlichkeit zu hierarchischen Bäumen besteht. Diese Ähnlichkeit kann den Anwendern bei der Erfassung des Dargestelltem helfen. Zudem hat die Kreisdarstellung geometrische Ähnlichkeit mit einer Medaille. Nach der Theorie der grafischen Semiologie \cite{bertin} verbessert diese Assoziation die Expressivität und Effektivität der Visualisierung. 
 
% Diskutieren sie die Expressivität und die Effektivität
Die Expressivität der Visualisierung ist außerdem durch die  Anforderung, die prozentualen Anteile im richtigen Verhältnis abzubilden, gegeben. Die konkreten Werte müssen nicht eingeblendet werden, um die entsprechenden Prozentwerte zu kennen. Da das SunBurst-Diagramm das mentale Modell von Anteilen direkt umsetzt, ist die Wahrnehmbarkeit und das Verständnis der Darstellung ebenfalls gegeben, sodass diese Visualisierungstechnik eine hohe Effektivität hat.

\subsubsection{Visualisierung Zwei}
% Einleitung
Die zweite Anwendungsaufgabe beschäftigt sich mit alternativen Länderrankings. Die Visualisierung soll diese Länderrankings darstellen und bei dem Vergleich der Platzierungen eines Landes unterstützen. Aus diesem Grund wurde ein Paralleles Koordinatensystem gewählt, beidem die Länderrankings an den Koordinatenachsen abgebildet werden.

% Präsentation der visuellen Abbildung und Kodierung der Daten
Bei der Formulierung der Anforderungen wurde bereits der Kompromiss eingegangen, keine Tabellen mit Platzierung und berechneten Werten darzustellen, sondern lediglich die Platzierung zu nutzen, da diese auch später als Vergleichskriterium dient. Somit stellt jede Achse in dem Parallelen Koordinatensystem ein Länderranking dar. Die Platzierungen werden in gleichen Abständen abgetragen, wobei sich der erste Platz an oberster Stelle befindet. Weiterhin werden die Platzierungen eines Landes auf den verschiedenen Achsen durch eine Linie zwischen den Koordinatenachsen miteinander verbunden. Diese wird für jedes Land mit einer anderen Farbe gezeichnet. Zusätzlich können die Achsen beliebig per \textit{Drag-And-Drop} angeordnet werden. Durch Anklicken einer Linie wird das entsprechende Land fokussiert und alle anderen im Hintergrund blass angezeigt.

% Begründung, warum und wie gut Desingentscheidung die Anforderungen erfüllt
Durch die Abbildung der Platzierungen auf die Koordinatenachsen wird die Anforderung erfüllt, dass die Darstellung der qualitativen Einflussfaktoren ebenfalls in einem Länderranking resultiert. Die Ähnlichkeit zu dem Medaillenspiegel ist für die Anwender leicht zu erkennen. Die Vergleichbarkeit wird durch die Linien zwischen den Koordinatenachsen erfüllt. Die Platzierungen werden nicht mehr explizit dargestellt, sondern müssen relativ zu der Beschriftung der Achse bestimmt werden. Jedoch kann durch den Winkel der Linie nun deutlich einfacher eine Verbesserung bzw. Verschlechterung in der Platzierung ermittelt werden. Durch die beliebige Anordnung der Achsen, kann der Nutzer alle Rankings leicht miteinander verglichen. In der Diskussion übe das mentale Modell des Rankings wurde festgestellt, dass dieses für die Anwendungsaufgabe nicht unbedingt benötigt wird. Es bietet lediglich einen guten Überblick. Deshalb wird per Klick auf eine Linie ein konkretes Land hervorgehoben und das Länderranking ausgeblendet.

% Begründung, warum gewählte visuelle Kodierung der Daten für Anwendungsaufgabe passend ist (z.B. Wahrnehmungsprinzip, Theorie über InfoVis)
Der Vergleich der Platzierungen wäre alternativ auch mit Datenpunkten in einem Scatterplot möglich. Dabei würden an der X-Achse die verschiedenen Rankings dargestellt werden und an der Y-Achse die entsprechende Platzierung. Diese Visualisierung wäre ähnlich zu dem parallelen Koordinatensystem, jedoch wird durch die Anordnung der Rankings entlang der X-Achse eine logische Abfolge suggeriert, die es nicht gibt. Vielmehr sollen die Nutzer die Anordnung der Rankings vertauschen, um direkte Vergleiche durchzuführen. Somit hat der Scatterplot keine Vorteile gegenüber den Parallelen Koordinatensystem. In der Diskussion über das mentale Modell des tabellarischen Rankings wurde festgestellt, dass dieses für die Visualisierung nicht unbedingt benötigt wird. Wir greifen es mit dieser Visualisierung jedoch bewusst auf, da die Nutzer solche Visualisierungen aus ihrem Alltag kennen und dadurch den dargestellten Zusammenhang schneller erfassen können. Eine Alternative wären hier zum Beispiel die Abbildung auf Farbe, wie bei einer HeatMap. Dabei können große Veränderungen in der Platzierung durch große farbliche Unterschiede erkenntlich gemacht werden. Dabei ist jedoch die Richtung der Veränderung (höhere/ niedrigere Platzierung) nicht ohne weitere Erklärungen klar. Ähnliche Probleme treten bei anderen alternativen Darstellungsformen mehrdimensionaler Daten, wie den Chernoff-Gesichter, auf. Die verschiedenen Rankings könnten durch unterschiedliche Attribute in dem Gesicht abgebildet werden. Beispielsweise würden zusammengekniffene Augen eine schlechte Platzierung in einem der Länderrankings bedeuten. Dadurch wäre der Vergleich zwischen den Attributen, welche die verschiedenen Platzierungen abbilden, aber sehr schwierig und von ungeübten Nutzern nicht zu erwarten. Daher ist in unseren Augen das Ranking immer noch die beste Darstellungsart und die Visualisierungstechnik der Parallelen Koordinaten optimal, um den Vergleich zwischen den Rankings durchzuführen.

% Diskutieren sie die Expressivität und die Effektivität
Die Platzierungen der Länder können mithilfe der Rankings korrekt dargestellt werden, sodass die Expressivität gegeben ist. Jedoch kann die Platzierung nicht genau abgelesen werden, was die Effektivität etwas einschränkt.

\subsubsection{Visualisierung Drei}
% Einleitung
Zur Beantwortung der Frage nach der sportlich erfolgreichsten Nation sollen die Anwender sich auch mit dem historischen Medaillenspiegeln auseinandersetzen. Dabei sollen sie Nationen mit regelmäßigen Medaillengewinnern und aufstrebende Nationen identifizieren. Die dazu benötigten Daten sind deutlich umfangreicher, als in den vorangegangenen Analysen. Die Visualisierungstechnik HeatMap unterstützt die Anwender beim Bearbeiten der Aufgabe durch eine kompakte Darstellung der Daten und ermöglicht die Analyse von Trends.

% Präsentation der visuellen Abbildung und Kodierung der Daten
Die Anzahl gewonnener Medaillen werden zunächst durch eine Farbe kodiert und als Zellen in der HeatMap dargestellt. Auf den Achsen der HeatMap werden die Länder und die Olympia-Jahre dargestellt. Die Jahre 1916, 1940 und 1944 wurden weg gelassen, da in diesen Jahren keine Olympiade stattfand. Wenn ein Land in einem Jahr an einer Olympiade nicht teilgenommen hat, wird der Wert Null kodiert. Für die Kodierung der Farben werden diskrete Stufen verwendet, welche für kleine Werte feingranularer sind, sodass Unterschiede zwischen kleinen Werten in wesentlichen Farbänderungen resultieren. Gleichzeitig werden alle Werte über 50 durch die gleiche Farbe kodiert, da mehr als 50 Medaillen nur die wenigsten Nationen gewonnen haben. Eine Legende an der Seite der Visualisierung ermöglicht die schnelle Interpretation der Farben. Durch \textit{Mouse-Hover} werden die konkreten Werte eingeblendet, sowie das entsprechende Land und Jahr an den Achsen hervorgehoben. Die Anwender können die Länder nach der Platzierung im Medaillenspiegel der Olympischen Sommerspiele 2024 oder nach dem sogenannten \enquote{All-Time} Medaillenspiegel sortieren. Der \enquote{All-Time} Medaillenspiegel ergibt sich aus der Anzahl aller jemals gewonnenen Medaillen einer Nation bei den Olympischen Sommerspielen.

% Begründung, warum und wie gut Desingentscheidung die Anforderungen erfüllt
Damit die Visualisierung die Anwender bei der Bearbeitung der Anwendungsaufgabe unterstützt, wurde eine effektive Kodierung der Daten gefordert. Darum wurden die Jahre ohne Olympische Spiele weg gelassen. Zudem werden die historischen Entwicklungen aller Länder in der HeatMap vereint, sodass die Anwender einen guten Überblick über alle Länder erhalten. Die Kodierung der Werte auf Farben steigert ebenfalls die Effektivität der Visualisierung, da farbliche Unterschiede schneller wahrgenommen werden als textuelle. Diese Eigenschaft beruht auf dem Konzept der \textit{Preattentive Recognition}. Demnach können bestimmte visuelle Eigenschaften, wie zum Beispiel Farbe, eine schnelle und unterbewusste Verarbeitung ermöglichen \cite{preAttent}. Mithilfe der Legende kann sichergestellt werden, dass alle Anwender die Abbildung der Werte auf Farben gleich interpretieren. Für die Vergleichbarkeit der Daten wurde zudem gefordert, dass sich die Farben für kleine Werte stark unterscheiden soll. Dies konnte durch die manuell festgelegten Farbstufen erreicht werden.

% Begründung, warum gewählte visuelle Kodierung der Daten für Anwendungsaufgabe passend ist (z.B. Wahrnehmungsprinzip, Theorie über InfoVis)
Die Anzahl gewonnener Medaillen sind quantitative Daten, welche sich nach \cite{bertin} am effektivsten auf Positionen abbilden lassen. Dies spricht für die Darstellung mithilfe eines Scatterplot, bei dem die Anzahl gewonnener Medaillen eines Landes pro Jahr als ein Datenpunkt visualisiert wird. Wenn diese Datenpunkte zusätzlich mit einer Linie verbunden werden ist auch eine Trendanalyse sehr einfach möglich. Der Vorteil der HeatMap gegenüber Scatterplots ist aber die Übersichtlichkeit bei der Darstellung vieler Länder. Im Scatterplot würde die Darstellung mehrerer Länder zu vielen überlappenden Linien führen, sodass eine Analyse nur schwierig möglich wäre. In der HeatMap ist die Vergleichbarkeit dagegen sehr gut möglich, da alle Länder gleichzeitig untereinander dargestellt werden können und dadurch ein direkter visueller Vergleich möglich ist. Zudem können in der HeatMap durch die Farbkodierung Muster in den Daten effizient erkannt werden. Ein Land, das regelmäßig viele Medaillen gewinnt wird durch eine durchgehend dunkle Zeile in der HeatMap visualisiert. Länder, die immer mehr Medaillen gewinnen werden durch eine immer dunkler werden Zeile dargestellt. Da diese Zellen direkt nebeneinander angeordnet sind, werden sie als Gruppe schnell wahrgenommen. Diese Wahrnehmung basiert auf der Feature Integration Theorie, wonach Begrenzungen (\textit{boundary detection}) durch Farben schnell wahrgenommen werden \cite{preAttent}. Die visuelle Struktur der HeatMap unterstützt die Trend-Analysen damit optimal.

\subsection{Interaktion}
% --> Die präsentierten Visualisierungstechniken müssen interaktiv zu einer Anwendung verknüpft werden. Die Interaktion mit einer Visualisierung soll in den anderen Visualisierungen zu einer Änderung führen.
% [ ] Erklären sie die möglichen Interaktionen mit den einzelnen Visualisierungen und die möglichen Verknüpfungen zwischen ihnen.
% [ ] Begründen Sie warum die konkreten Interaktionen umgesetzt wurden und welche Zwecke für die Anwenderinnen mit ihnen unterstützt werden.
% [ ] Begründen sie ebenfalls warum sie andere Interaktionsmöglichkeiten nicht umgesetzt haben.
% --> Wenn sie keine der geforderten Interaktionen umsetzen, erhalten Sie im gesamten Projekt deutlichen Punktabzug.
Mithilfe dieser Visual-Analytics-Anwendung sollen die Anwender den Medaillenspiegel untersuchen und hinterfragen. Die Analyse beginnen sie bei dem für sie bekannten Medaillenspiegel und gelangen durch Interaktionen zu den verschiedenen Visualisierungen, welche sie bei der Bearbeitung der Anwendungsaufgaben unterstützen.

Da die gleichzeitige Analyse aller 204 teilnehmenden Länder nicht möglich ist, sollten die Anwender ein konkretes Land auswählen können, welches sie dann im SunBurst-Diagramm und den verschiedenen Länderranking untersuchen und analysieren können. Die Anwendungsaufgabe der ersten Visualisierung ist die, im Medaillenspiegel dargestellte Anzahl gewonnener Medaillen zu hinterfragen und sich mit deren Verteilung nach Sportarten und Disziplinen auseinanderzusetzen. Daher wurde eine Interaktion implementiert, bei der durch Klick auf ein Land im Medaillenspiegel, das entsprechende SunBurst-Diagramm geladen wird. Die Anwender sollen die Analyse des ausgewählten Landes dann gleich fortsetzen, weshalb dieses in der Visualisierung der verschiedenen Länderrankings ebenfalls hervorgehoben ist.

Für den Vergleich der verschiedenen Länderrankings sollten die Koordinatenachsen beliebig angeordnet werden können, um die Platzierungen zwischen zwei verschiedenen Länderrankings direkt vergleichen zu können. Da auch diese Interaktion bei der Bearbeitung der Anwendungsaufgabe hilft, wurde diese ebenfalls implementiert. Die Koordinatenachsen können per \textit{Drag-And-Drop} verschoben werden. Damit die einzelnen Länderrankings auch im Ganzen analysiert werden können, kann per Klick auf eine Koordinatenachse der Medaillenspiegel durch dieses alternative Länderranking ersetzt werden.

In der ersten Anwendungsaufgabe sollen die Anwender sich mit der Verteilung der Medaillen nach Sportarten und Disziplinen auseinandersetzen. Dabei betrachten sie die gewonnenen Medaillen der Olympischen Sommerspiele 2024. Um auch die anderen Jahre analysieren zu können ist eine Interaktion zwischen der HeatMap und SunBurst-Diagramm nötig, bei der auch für vergangene Jahre die Medaillenverteilung geladen werden kann. Wir haben diese Interaktion nicht implementiert, da die Menge der Sportarten bei jeder Olympiade anders ist und ein Vergleich zwischen den Jahren nicht sinnvoll erscheint.